\documentclass[11pt]{article}
\usepackage[utf8]{inputenc}
\usepackage{amsmath,amssymb}
\usepackage[margin=1in]{geometry}
\usepackage{hyperref}
\emergencystretch=3em

\title{Leading asymptotic of the standard integral for polynomial data:
only $\xi(0)$ and $\eta(0)$ survive}
\author{Claude, with Daniel Murfet}
\date{2026-09-06}

\begin{document}
\maketitle
\sloppy

\begin{abstract}
The capstone of the one-dimensional programme: for polynomial data
$\xi(u)=\xi_0+\sum_{j<D}c_ju^{j+1}$ and $\eta(u)=\sum_{i\le E}e_iu^i$ with $\eta(0)\ne0$,
\[
\int_0^b u^h\eta(u)\,e^{-\beta nu^{2k}+\beta\sqrt n u^k\xi(u)}\,du
\;\sim\;\eta(0)\,\tfrac1{2k}\,S_{(h+1)/(2k)}(\xi(0))\,n^{-(h+1)/(2k)}\qquad(n\to\infty),
\]
a genuine \texttt{Asymptotics.IsEquivalent}. This is the leading term of
\texttt{cor:standardintegralexp} in one dimension for general (polynomial) data: the leading exponent
is the first element of $\Lambda(h,k)$ and only the chart-origin values $\xi(0),\eta(0)$ enter the
coefficient. En route, the polynomial-$\xi$ two-term bound of unit~32 is packaged as
$Z-\tfrac1{2k}n^{-(h+1)/(2k)}S(\xi_0)=O(n^{-(h+2)/(2k)})$. Zero \texttt{sorry}/\texttt{axiom}.
\end{abstract}

\section{Setting}
Units~27--32 establish the expansion for $\eta\equiv1$ and for monomial or polynomial $\xi$. The
paper's standard integral carries a weight $\eta$ (Jacobian and observable factors); by linearity a
polynomial $\eta$ redistributes the problem over shifted powers $u^{h+i}$, each contributing at order
$n^{-(h+i+1)/(2k)}$. So the $i=0$ term dominates and the rest is $O(n^{-(h+2)/(2k)})$, and the
polynomial-$\xi$ correction is of the same order.

\section{The things formalised}
{\small
\begin{verbatim}
theorem isLittleO_rpow_neg_rpow_neg {a b : ℝ} (hab : b < a) :
    (fun n : ℝ => n ^ (-a)) =o[atTop] fun n : ℝ => n ^ (-b)
theorem isBigO_rpow_neg_rpow_neg {a b : ℝ} (hab : b ≤ a) :
    (fun n : ℝ => n ^ (-a)) =O[atTop] fun n : ℝ => n ^ (-b)

theorem standardIntegral1D_polynomial_leading_isBigO ... :
    (fun n => Z_chart(ξ) n - (1/(2k)) n^(-(h+1)/(2k)) S_{(h+1)/(2k)}(ξ₀))
      =O[atTop] fun n => n ^ (-((h:ℝ)+2)/(2k))

theorem standardIntegral1D_weighted_leading_isEquivalent ... (he0 : e 0 ≠ 0) :
    (fun n => ∫ u in Ioc 0 b, u^h * (∑ i ∈ range (E+1), e i * u^i) * exp(...ξ(u)))
      ~[atTop] fun n => e 0 * ((1/(2k)) * S_{(h+1)/(2k)}(ξ₀)) * n^(-(h+1)/(2k))
\end{verbatim}
}

\section{Proof strategy}
The power helpers are \texttt{isLittleO\_of\_tendsto'} with \texttt{tendsto\_rpow\_neg\_atTop} (ratio
$n^{-(a-b)}\to0$) and \texttt{IsBigO.of\_bound 1} with \texttt{rpow\_le\_rpow\_of\_exponent\_le}. The
polynomial-$\xi$ big-$O$ is unit~29's packaging applied to unit~32: the $D$ first-order terms and the
remainder are each dominated by $n^{-(h+2)/(2k)}$ for $n\ge1$ (exponent comparison), the tails by the
exponential scales. The capstone splits $\eta$ by \texttt{Finset.sum\_range\_succ'} into $e_0$ and the
rest, rewrites the weighted integral by linearity (\texttt{Finset.mul\_sum}, \texttt{integral\_finsetSum}),
and shows the difference from $e_0A_0n^{-\mu_1}$ is $O(n^{-\nu})$ with $\nu=(h{+}2)/(2k)$: the $e_0$ part
by the big-$O$ at $h$, each $i\ge1$ part by the big-$O$ at $h+i$ plus its own leading term, all pulled down
to $n^{-\nu}$ by the power helper and summed with \texttt{IsBigO.sum}. Then $n^{-\nu}=o(n^{-\mu_1})$
and \texttt{IsLittleO.const\_mul\_right} with the nonzero coefficient $e_0A_0$ give \texttt{IsEquivalent}.

\section{Roadblocks and resolutions}
\begin{itemize}
\item \texttt{gcongr} on $\tfrac{h+2}{2k}\le\tfrac{h+2+1}{2k}$ paired the summands as $2$ against $1$
and left `$2\le1$'; for exponent comparisons use \texttt{div\_le\_div\_iff\_of\_pos\_right} then
\texttt{linarith}. Likewise \texttt{gcongr} on $n^{-X}\le n^{-\mu}$ reduces to the \emph{exponent}
inequality, not the power inequality one has in hand --- pass explicit \texttt{mul\_le\_mul} chains.
\item A \texttt{calc} step with `\_ + \_' placeholders in the middle expression was unified against the
wrong subterms; spell the intermediate expression out (or name the bound as a separate \texttt{have}).
\item \texttt{(IsBigO.add ?\_ ?\_).congr\_left \dots} left a typeclass problem stuck because the two
functions were undetermined; name both big-$O$ facts first. \texttt{IsBigO.sum} produces a Pi-type
\texttt{Finset.sum} of functions, converted to a pointwise sum by \texttt{Finset.sum\_apply} inside
\texttt{congr\_left}.
\item The final \texttt{congr'} with \texttt{EventuallyEq.rfl} also got stuck; \texttt{IsLittleO.congr\_left}
with the linearity identity is the robust route.
\end{itemize}

\section{What was Mathlib, what was new}
Mathlib: \texttt{tendsto\_rpow\_neg\_atTop}, \texttt{isLittleO\_of\_tendsto'},
\texttt{rpow\_le\_rpow\_of\_exponent\_le}, \texttt{IsBigO.sum}, \texttt{IsBigO.const\_mul\_left},
\texttt{IsLittleO.const\_mul\_right}, \texttt{Finset.sum\_range\_succ'}, \texttt{Finset.sum\_apply}. New:
the two power-comparison helpers and the leading asymptotic for polynomial $(\xi,\eta)$.

\section{Lessons}
\texttt{gcongr} is excellent when the shapes on both sides match term-for-term and treacherous
otherwise --- with sums of different lengths or when a monotone function (here \texttt{rpow} in the
exponent) rewrites the goal into a form different from the hypothesis one holds. In asymptotic
bookkeeping, name every big-$O$ fact with its function written out before combining.

\section{Outcome}
The one-dimensional analytic content of grammar \S4.2 is now formalised for polynomial data: exact
expansions, explicit remainders, and asymptotic statements at leading, second and all orders (units
22--33). The remaining gaps are structural rather than analytic: general analytic $\xi,\eta$ (limits of
the polynomial statements), the multivariate tree with its $\log n$ powers, and \S4.3's stochastic
layer.
\end{document}
